\documentclass[12pt,a4paper]{article}
\usepackage[utf8]{inputenc}
\usepackage[T1]{fontenc}
\usepackage[spanish,es-noquoting,es-noshorthands]{babel}
\usepackage{geometry}
\geometry{a4paper, margin=2.5cm}
\usepackage{xcolor}

\title{\textbf{Código de Ética y Reglas de Uso}\\ \large Suite de Auditoría de Seguridad Informática}
\author{Prof. César Rodríguez \\ \small Curso: Seguridad Informática}
\date{\today}

\begin{document}

\maketitle

\section*{Preámbulo}
Las herramientas y módulos desarrollados durante el curso de Seguridad Informática, incluyendo escaneo de puertos, enumeración, recolección de información (OSINT) y ataques de fuerza bruta, tienen un potencial dual. Si bien son fundamentales para comprender cómo proteger sistemas, también pueden ser empleadas de forma maliciosa para comprometer la infraestructura de terceros.

Por consiguiente, el desarrollo y uso de este software requiere un estricto apego a los principios éticos y legales que rigen la profesión de la ciberseguridad.

\section*{1. Declaración de Responsabilidad}
El uso de cualquier script, módulo o técnica descrita en este repositorio es responsabilidad exclusiva del estudiante. Ni el profesor César Rodríguez, ni la institución académica asumen responsabilidad alguna por el mal uso directo o indirecto del código aquí desarrollado.

\section*{2. Reglas de Uso Obligatorias}
Al descargar, ejecutar o modificar el código fuente del orquestador (\texttt{auditoria.py}) y sus módulos asociados, el estudiante acepta cumplir las siguientes reglas ineludibles:
\begin{itemize}
    \item \textbf{Autorización Explícita:} Queda terminantemente prohibido ejecutar escaneos, enumeraciones o ataques de diccionario contra cualquier servidor, dominio, red o dirección IP sin el consentimiento previo, explícito y por escrito de su propietario.
    \item \textbf{Entornos de Prueba:} Las pruebas de los módulos deben realizarse de manera estrictamente local (\texttt{127.0.0.1}), en máquinas virtuales propias, o contra servidores de prueba autorizados expresamente para fines educativos (por ejemplo, \texttt{scanme.nmap.org}).
    \item \textbf{Impacto Cero:} Se prohíbe el uso de estas herramientas para realizar ataques de Denegación de Servicio (DoS), interrupción de operaciones, o extracción no autorizada de datos.
\end{itemize}

\section*{3. Consecuencias Académicas y Legales}
El incumplimiento de este Código de Ética será tratado como una falta disciplinaria muy grave. Dependiendo de la severidad de la infracción, puede conllevar desde la reprobación automática del curso y expulsión de la institución, hasta el reporte formal a las autoridades competentes, dado que el acceso no autorizado a sistemas informáticos es un delito penal.

\vspace{1.5cm}

\begin{center}
\fbox{
\parbox{0.9\textwidth}{
    \vspace{0.3cm}
    \begin{center}
        \textbf{Acuerdo de Compromiso}
    \end{center}
    \vspace{0.2cm}
    
    \textit{He leído y comprendido íntegramente este Código de Ética. Me comprometo a utilizar las herramientas desarrolladas en las Fases I, II y III exclusivamente para fines académicos, defensivos y éticos.}
    
    \vspace{1.5cm}
    Nombre del Estudiante: \hrulefill \\
    
    \vspace{1cm}
    Firma del Estudiante: \hrulefill \\
    
    \vspace{1cm}
    Fecha: \hrulefill
    \vspace{0.3cm}
}
}
\end{center}

\end{document}